\subsubsection{Alternative-Phase Deviation Tests}
\label{sec:phase-hardware}
This experiment evaluates $Q_N^\star$ separately from the original
phase. At $N=3$, $Q_3^\star=Q_3$, so that size retests the original phase
with Dove deviations added. On 7 September 2026, a preregistered plan fixed a
seven-qubit chain, $[141,142,143,144,145,146,147]$, on ibm\_fez and used
its prefixes for $N=3,\ldots,7$. An 18-circuit pilot run preceded the
complete 70-circuit batch. The complete batch contains a cooperative profile at
each size, both Hawk and Dove deviations for every player, original-phase
controls, an unrestricted counterstrategy, and two readout calibrations.
Every circuit uses 4096 shots. Each circuit is checked against its dense
operator and its compiled ideal output, and the batch is rehearsed on a
device noise model before submission.

The primary analysis uses raw implemented payoffs: calibration circuits
are diagnostic and no readout correction or ZNE is applied to these tests.
With $M=50$ primary deviations and $\alpha=0.05$, the simultaneous lower
bound on each cooperative-minus-deviation gap is
\begin{equation}
g_j^{\rm L}(a)=\widehat g_j(a)
-\sum_{s\in\{\mathrm{coop},\mathrm{dev}\}}
V\sqrt{\frac{\log(2M/\alpha)}{2n_s}},
\end{equation}
where $n_s=4096$. Hoeffding's inequality and a union bound cover the two
one-sided mean errors for every gap. These bounds assume independent
shots from stationary circuit distributions; they do not cover temporal
drift. The preregistered approximate-equilibrium tolerance is
$-g_j^{\rm L}(a)/(V/N)\le0.05$ for every deviation of every player.
The overall test requires all five sizes to pass.

In the first complete batch, every measured gap is positive, but the
simultaneous criterion passes only at $N=3,4,5,6$. The minimum observed
gaps are $0.2360$, $0.5677$, $0.2157$, $0.2627$, and $0.0109$,
respectively. The corresponding maximum upper bounds on normalized
deviation gain are $0.0058$, $-0.3240$, $0.0349$, $-0.0285$, and
$0.4074$. Thus the overall test fails: a positive point gap
at $N=7$ does not establish the requested tolerance. At $N=4,6$ the
simultaneous lower gaps are strictly positive; at $N=3,5$ the supported
statement is approximate equilibrium at the declared tolerance.

This all-player test locates an implementation boundary absent from mean
welfare alone. At $N=7$, target population is $0.7495$ and analytic-baseline
retention is $0.9600$, below the $0.9922$ that uniformly random outcomes
would retain, and incentive compatibility is unresolved. The unrestricted
counterstrategy, run as a diagnostic, gains between $2.4604$ and $2.7052$
over cooperation for player zero across the five sizes, consistent in
direction with the analytic counterexample.

A second complete batch reused the same circuits, shots and decision rule. Each batch is judged separately;
their counts are not pooled to rescue a failed criterion. The repeat
supports the criterion at $N=4,6$, but not at $N=3,5,7$
(Table~\ref{tab:phase-repeat}). All 50 point gaps are positive in the first
batch; one is slightly negative in the repeat, at $N=7$. This is not a
statistically established profitable deviation. Both batches fail the
overall five-size test, while the positive lower gaps at $N=4,6$ reproduce.
The executions occurred minutes apart on the same chain; this is a
within-session repetition, not evidence across independent calibration
epochs. The pilot and two complete batches consumed $21+77+77=175$ seconds
of provider-reported charged QPU time.

\begin{table}[t]
\centering
\caption{Alternative-phase all-player tests, judged separately in each
batch. $g_{\min}$ is the smallest observed payoff gap; $u_{\max}$ is the
largest simultaneous upper bound on deviation gain normalized by $V/N$.
The preregistered tolerance requires $u_{\max}\le0.05$.}
\label{tab:phase-repeat}
\begin{tabular}{crrrr}
\toprule
& \multicolumn{2}{c}{First batch} & \multicolumn{2}{c}{Repeat}\\
$N$ & $g_{\min}$ & $u_{\max}$ & $g_{\min}$ & $u_{\max}$\\
\midrule
3 & 0.2360 & 0.0058 & 0.1362 & 0.0806\\
4 & 0.5677 & -0.3240 & 0.5716 & -0.3279\\
5 & 0.2157 & 0.0349 & 0.1861 & 0.0719\\
6 & 0.2627 & -0.0285 & 0.2483 & -0.0069\\
7 & 0.0109 & 0.4074 & -0.0033 & 0.4322\\
\bottomrule
\end{tabular}
\end{table}
